\begin{lemma}[equivariant retraction of the cusp body]\label{push:src:C.3.8}
For a sufficiently small closed cusp disc \(\overline\Delta\),
the compact manifold with boundary
\(V=f^{-1}(\overline\Delta)\) has a \(\rho\)-equivariant strong
deformation retraction onto \(W=f^{-1}(0)\). Consequently
\(V/\rho\) strongly deformation retracts onto \(H=W/\rho\),
and the fixed sphere \(F\) is primitive and nonzero in
\(H_2(V/\rho;\mathbb Z)\). It cannot bound a three-chain,
and hence cannot bound a three-ball, inside this cusp orbit body.
\end{lemma}
\Uses{the hypotheses, notation and definitions stated here; \cref{push:src:C.3.7}}

The existence of the torus-equivariant strong deformation retraction is already \citet[Proposition~7.2, p.~41; proof completed on p.~50]{AC26}. Restricting that equivariance to the scalar circle gives the retraction used here. The additional assertion is the primitive nonzero class of the fixed sphere in the quotient body; its identification uses the quotient calculation of \cref{push:src:C.3.7}. The complete retraction argument below is retained.

\begin{proof}
Write \(q\) for the cusp coordinate and set \(g=|q\circ f|^2\).
Choose a smooth invariant Hermitian metric by averaging any smooth
Hermitian metric over the compact circle. The negative gradient
flow of \(g\) is equivariant, fixes \(W\) pointwise, and points
strictly inward on \(\partial V\). Its flow \(\varphi_t\) exists
for every \(t\geq0\), since \(V\) is compact.

We prove the uniform convergence to \(W\) that is needed for a
strong deformation retraction. At every point of \(W\), the
actual semistable cusp charts give holomorphic coordinates with
\(q\circ f=z_1\cdots z_k\), where \(1\leq k\leq3\).
An initial holomorphic unit can be absorbed in one coordinate.
For a point off \(W\), some \(i\) satisfies
\( |z_i|\leq|q\circ f|^{1/k}\), and therefore
\[
 \left|\frac{\partial(q\circ f)}{\partial z_i}\right|
       \geq |q\circ f|^{1-1/k}.
\]
Uniform comparison with the chosen metric on a finite cover of
the compact fibre, and \(g\leq1\) after shrinking, give a
constant \(c>0\) such that
\[
 |\nabla g|\geq c g^{5/6}
\]
near \(W\). On the compact complement of a smaller such
neighbourhood, \(g>0\) and \(\nabla g\ne0\), because
\(f\) is a submersion off its central fibre. Decreasing \(c\)
extends this inequality to all of \(V\setminus W\).

Along a trajectory put \(h(t)=g(\varphi_t(x))\). Then
\(h'=-|\nabla g|^2\leq-c^2h^{5/3}\). If
\(G=\max_V g>0\), integration gives the uniform bound
\[
 h(t)\leq\bigl(G^{-2/3}+(2c^2/3)t\bigr)^{-3/2}.
\]
The remaining length of the trajectory starting at time \(t\)
is bounded by
\[
 \int_t^\infty|\nabla g(\varphi_s(x))|\,ds
 =\int_0^{h(t)}\frac{dv}{|\nabla g|}
 \leq\frac6c h(t)^{1/6}.
\]
One may first integrate to a finite time; the displayed uniform
bound and monotone convergence justify the endpoint expression.
Thus every trajectory has a limit in \(W\), and the tail-length
bound tends uniformly to zero. The maps \(\varphi_t\) converge
uniformly to a continuous equivariant retraction \(r:V\to W\).
The homotopy given by \(\varphi_{s/(1-s)}\) for \(s<1\)
and by \(r\) at \(s=1\) is continuous, equivariant, and fixes
\(W\) throughout. This is the required strong deformation
retraction. It descends to the orbit quotient, so the homology
and primitivity statement transfers from \ref{push:src:C.3.7}. The same argument
does not supply a global Seifert surface in the full orbit sphere.
\end{proof}

